    \section{Non-Transformer Architectures and Run Protocols}
\label{app:wing-models}

\subsection{Evidence Supporting the Four-Model Comparison}
\label{app:selection-rationale}

Unified reevaluation of earlier NEXT checkpoints provides retrospective
context for the representative architectures. In the alternative-model batch, the multi-view CNN
has matched AUC 0.955804, compared with 0.947782 for the multiscale CNN and
0.928273 for the dense 3D CNN; Static GINE leads the graph candidates in that
batch at 0.981008. In the subsequent candidate batch, BiGRU exceeds the
tested dilated TCN (0.933864 versus 0.926085). PointMamba-lite, the only recorded
state-space candidate, has matched AUC 0.925386 and $I=0.971387$; its inclusion broadens the
architectural comparison rather than following an overall accuracy ranking.

These are held-out test comparisons on 115,499 events from earlier
exploratory batches. They are distinct from the 116,549-event reevaluation
used for the available NEXT classic results in Table~\ref{tab:benchmark-main}.
PointMamba-lite is included in
Table~\ref{tab:next-classification}. The original exploratory batches used different evaluator-code fingerprints;
the numbers above now use the same regular-plus-overflow protocol as the
current classic results while retaining the earlier checkpoints and test set.
Their distinct test population is not pooled with the retained campaign. The four search groups are recorded in the model-family
presentation and are reproduced in Table~\ref{tab:next-classification}.
They organize the candidate inventory; the surviving records do not
document a prospective validation-based rule selecting one winner from
each group. The retained architectures span three of the four groups,
including two distinct mechanisms from the sequence/state-space group.
The structural rationale and early test comparisons therefore support the
comparison set without establishing a recovered formal selection rule.

\begin{table}[t]
\centering
\small
\setlength{\tabcolsep}{4pt}
\caption{Input adaptations for the four non-Transformer architectures. GINE, BiGRU, and PointMamba-lite share each dataset's point representation. Model dimensions and representation details are given in Appendix~\ref{app:wing-models}.}
\label{tab:wing-adaptations}
\begin{tabular}{p{0.14\textwidth}p{0.35\textwidth}p{0.43\textwidth}}
\toprule
Dataset & MV-CNN input & Graph and sequence input \\
\midrule
NEXT & Three normalized spatial energy projections & Centered 15\,mm voxels; energy fraction and log hit count \\
SuperNEMO & Three binary tracker projections & Centered 15\,mm tracker cells; occupancy fraction and log hit count \\
MJD & Three-resolution waveform STFT & 512 pooled waveform segments with mean and RMS features \\
EXO-200 & STFT of serialized $226\!\times\!300$ waveform & Same waveform construction after channel-major serialization \\
\bottomrule
\end{tabular}
\end{table}


\paragraph{Parameter-count provenance.}
Parameter counts in Table~\ref{tab:next-classification} are the saved
counts of trainable neural parameters, checked against the corresponding
run summaries and reevaluated metrics. Neural parameter count is
inapplicable to XGBoost. The earlier NEXT PointMamba-lite run summary records
316,993 trainable parameters, consistent with \texttt{Benchmarking table.xlsx};
this is displayed as 0.3170\,M in Table~\ref{tab:next-classification} and
317\,K in Table~\ref{tab:benchmark-main}. Its retained-campaign predictions
remain unavailable. Figure~\ref{fig:next-capacity} includes only the
retained-campaign configurations with available predictions.

\paragraph{Spatial representations.}
NEXT MV-CNN uses three $128\times128$ projections with 30\,mm bins. For point-based models, hits are centered by the complete-event energy centroid and merged into 15\,mm cells. Voxel centers are recentered and divided by 1,000\,mm; the two features are the voxel's fraction of complete-event energy and $\log(1+n)$, where $n$ is its hit count. Events exceeding 512 occupied cells retain the highest-energy cells with deterministic tie breaking, without renormalizing the retained energy fractions. SuperNEMO instead uses tracker coordinates: its CNN projects binary occupancy on $128\times128$ grids with 44\,mm bins. Its point adapter centers the tracker hits geometrically and merges 15\,mm cells, replacing the unavailable hit-energy fractions with occupancy fractions. It raises an error above 512 cells instead of truncating. The non-Transformer SuperNEMO inputs exclude $E_1$, $E_2$, and tracker drift radii. These input choices restrict direct energy access but do not establish independence from energy.

\paragraph{Waveform representations.}
MJD subtracts the mean of the first 200 waveform samples and divides by the maximum absolute event amplitude when nonzero. EXO-200 subtracts a baseline separately for each channel using its first 200 time samples, divides by the event-wide maximum absolute amplitude, and serializes the $226\times300$ array in channel-major order. For these datasets, MV-CNN uses Hann-window STFTs with FFT sizes $(64,128,256)$ and hops $(16,32,64)$, followed by a $\log(1+|\mathrm{STFT}|)$ transform and resizing to $128\times128$. The point adapter uses adaptive pooling to form 512 segment means and root-mean-square amplitudes. Its three coordinates are normalized time, amplitude-normalized segment mean, and the first difference of that normalized mean; its two features are the segment mean and RMS. These are pulse-shape coordinates, not calibrated detector positions.

\paragraph{Backbone details.}
MV-CNN has four residual stages with two blocks each, base width 16, and a 256-dimensional fusion layer. GINE uses hidden width 128, five message-passing layers, $k=12$ neighbors, and a 160-dimensional classifier. BiGRU uses a 96-dimensional input embedding, two bidirectional layers with hidden size 128, and a 256-dimensional classifier. The two Hilbert orders quantize each event's point bounding box to 10 bits per axis; the second swaps the first two axes before applying the same ordering. PointMamba-lite uses three blocks with model width 128, inner width 192, state dimension 16, and 32-position scan chunks. Its guarded PyTorch scan is not the official fused Mamba implementation. The backbones use dropout 0.1. MJD GINE has four output logits $z_{ij}$, one per reference label, and uses
\begin{equation}
 q_i=\prod_{j=1}^{4}\sigma(z_{ij}),\qquad
 s_i=\log q_i-\log(1-q_i),
\end{equation}
where $\sigma$ is the logistic sigmoid and $s_i$ is its evaluated clean-event score. This product is the implemented score construction; it is not a demonstrated factorization of the labels' joint conditional distribution. The computation uses summed log-sigmoids and clamps the log probability below zero for numerical stability. EXO-200 uses the same GINE backbone with one output.

\begin{table}[t]
\centering
\small
\caption{Specified common protocol versus recorded training settings for the audited non-Transformer models. Every listed configuration has an upper limit of 50 epochs, weight decay $10^{-4}$, gradient-norm limit 1, and seed 42. NEXT settings refer to the three available retained-campaign runs. ``Other datasets'' means EXO-200, MJD, and SuperNEMO.}
\label{tab:wing-training}
\begin{tabular}{llrrr}
\toprule
Scope & Architecture & Batch & Initial learning rate & Patience \\
\midrule
Specified protocol & Neural models & 64 & $5\times10^{-4}$ & 5 \\
NEXT recorded & MV-CNN, GINE, BiGRU & 64 & $5\times10^{-4}$ & 5 \\
Other datasets & MV-CNN & 16 & $10^{-3}$ & 12 \\
Other datasets & GINE, BiGRU & 16 & $5\times10^{-4}$ & 12 \\
Other datasets & PointMamba-lite & 4 & $3\times10^{-4}$ & 12 \\
\bottomrule
\end{tabular}
\end{table}


\paragraph{Run completion and randomness.}
The audited non-Transformer classification loaders retain the sampled class proportions;
their binary losses do not use class weights, and these paths contain no
synthetic data augmentation. NEXT shuffles training slices and buffered
events, waveform loaders shuffle training events, and SuperNEMO records
proportional sampling without replacement. Input normalization and ordering
are dataset-specific preprocessing rather than energy matching of training
events.

In the NEXT campaign reported here, MV-CNN, GINE, and BiGRU finish after 19, 36, and 26 epochs, selecting epochs 14, 31, and 21, respectively, by validation inclusive AUC. The loader batch size is 64; the wrapper evaluates these backbones in chunks of 16, concatenates their logits, and applies one batch loss and optimizer step. The corresponding run summaries record Python 3.11.15, PyTorch 2.11.0+cu128, an NVIDIA GeForce RTX 5090, and bfloat16 mixed precision. This environment evidence applies to those runs rather than automatically to every dataset. The audited non-Transformer training paths add neither class-weighted losses nor data augmentation; the current SuperNEMO stream retains the source class proportions without replacement.

For SuperNEMO, the workbook records 21, 50, 22, and 9 training epochs for MV-CNN, GINE, BiGRU, and PointMamba-lite, with selected epochs 16, 43, 17, and 4, respectively. It records 47 epochs and selected epoch 43 for MJD PointMamba-lite. Earlier logs were incomplete; the current results and epoch counts follow the workbook, whose evaluation status does not establish training completion. The audited non-Transformer configurations set seed 42 and disable enforced deterministic algorithms. Mixed precision is enabled in these configurations except for MJD BiGRU. Seeds therefore do not imply bitwise reproducibility.

\paragraph{Screening scope.}
The current NEXT registry contains three legacy CNNs and twenty alternative classifiers, including point, graph, sequence, state-space, sparse-convolution, hybrid, mixer, persistence, and boosted-tree approaches. A classification entry need not correspond to a completed retained-campaign run: PointMamba-lite has an earlier exploratory result but remains unavailable in that campaign. The search taxonomy groups the three legacy projection CNNs together with projection, patch-mixer, dense-voxel, and sparse-voxel alternatives. Point-set and graph methods form separate groups; ordered sequence methods, the state-space model, topological classifiers, and the CNN--GNN hybrid form the fourth. This representation-based organization differs from the finer model-family labels in the run registry. Checkpoint selection is documented within completed runs, while a prospective rule for choosing the cross-dataset architectures is not. The boosted-tree classifier in the extended NEXT inventory uses up to 500 estimators, learning rate 0.04, maximum depth 4, subsampling 0.85, column subsampling 0.9, early-stopping rounds 12, and random seed 42. Its estimator settings, rather than the neural epoch, optimizer, and clipping fields in the generic run configuration, determine its training procedure.
